\section{Stops for endpoints with small effective size}
\label{sec:stops}

The inverse and variance estimates leave compact endpoints whose large
prime factors all align to a rational of small effective size.  We now
stop the tree earlier on such branches.  At these stops we compare the
exact survivor counts with their reference subtrees after averaging over
one large prime.  All prime cutoffs at a nonempty node in this section
are strict: if its last prime is $P$, the local survivor count avoids
the forbidden classes at the primes $t<P$.

Fix the exponent $C_s$ supplied by \cref{cor:hard-endpoints}.  The
parameters of the regular boxes, including $C_{\rm len}$, are also
regarded as fixed.  The bin width $\xi$ will be chosen sufficiently small
only after the other fixed parameters of this section.  As before, let
$h\asymp\xi/\log w$ be the exponent width of a bin and put
\[
 \Delta=2C_{\rm len}\xi.
\]
The representative of a bin is its right endpoint in exponent
coordinates.  By \cref{prop:regular-boxes}, the movement of any prefix
gap within a box of length at most $C_{\rm len}\log B$ is at most
$\Delta$, for sufficiently large $z$.

\subsection{Owners, tags, and a stopping rule}

\begin{definition}\label{def:stops-owner}
A reduced rational $A/D$, with $D>0$, is \emph{eligible} if
\begin{equation}\label{eq:stops-eligible}
 D\le w^{C_s},\qquad |A|\le Yw^{C_s+2}.
\end{equation}
Fix a constant $\eta>0$, to be bounded by an absolute constant below.
A search bin has an \emph{owner} if some eligible rational aligns to
at least a $1-\eta$ fraction of its primes, by prime count.  The owner
is that rational.  Along a descending-prime path, its \emph{tag prime}
is the first search-bin factor that aligns to the owner of its bin,
if such a factor occurs.  Its \emph{tag rational} is this owner and
is retained for the rest of the path.
\end{definition}

The owner is unique for all sufficiently large $z$, provided
$\eta<1/2$.  Indeed, if $A_1/D_1\ne A_2/D_2$ are eligible, every
prime aligning to both divides the nonzero integer
$A_1D_2-A_2D_1$, whose absolute value is at most
\[
 2Yw^{2C_s+2}=z^{2+o(1)}.
\]
Consequently, in a bin with lower prime endpoint $R$, their common
aligning primes number at most
\begin{equation}\label{eq:stops-overlap}
 O\left(\frac{L}{\log R}+1\right).
\end{equation}
For a search bin, $\log R\ge w^{1/4}\log w$, while its prime count
is $\asymp\xi R/\log R$.  The bound in
\eqref{eq:stops-overlap} is a uniformly negligible fraction of this
count.  Two subsets of relative size at least $1-\eta$ would have an
intersection of relative size at least $1-2\eta$, proving uniqueness.
Also, every prime of a search bin eventually exceeds every eligible
denominator.

Choose fixed constants $b_0,b_1$ with $b_1>b_0$, taking $b_0$
sufficiently large compared with $C_s+1$.  Further conditions on these
constants will be stated below.  For a prefix with bin representatives
$x_1^{\rm rep},\ldots,x_j^{\rm rep}$, write
\[
 r_j^{\rm rep}=r_0-\sum_{i=1}^j x_i^{\rm rep}.
\]

\begin{definition}\label{def:stops-rule}
An included, nonempty, even node is a \emph{stop} if it is the first
node along its branch satisfying all the following conditions:
\begin{enumerate}[label=(\roman*)]
\item Its length is at most $C_{\rm len}\log B$.  Its tag has been
attained, the tag bin has multiplicity one in this prefix, and every
prime of the prefix aligns to its tag rational.
\item Every nonempty even prefix ending at position $j$ satisfies
\begin{equation}\label{eq:stops-safety}
 r_j^{\rm rep}\ge \mathcal H(x_j^{\rm rep})+3\Delta.
\end{equation}
\item At the present node, its last representative exponent and
representative ratio satisfy
\begin{equation}\label{eq:stops-window}
 b_0\le x_j^{\rm rep}\le b_1,
 \qquad
 2.06\le r_j^{\rm rep}/x_j^{\rm rep}\le2.16.
\end{equation}
\end{enumerate}
The empty root is not subjected to \eqref{eq:stops-safety}.
\end{definition}

This is a prefix-free choice of lower nodes.  The rule applies to the
entire included tree; the inverse and variance exceptions classify
endpoints but do not restrict this rule.
The representative safety
condition ensures that varying prime choices inside their bins, while
preserving multiplicities and distinctness, preserves admission of all
even prefixes by \eqref{eq:admission}.  One may see this directly from
the right-endpoint convention: the actual gap is at least the
representative gap and the actual last exponent is at most its
representative.  Alternatively, the movement bound and the
$2$-Lipschitz property of $\mathcal H$ give the same conclusion with room to
spare.  Odd children of lower nodes require no admission test.

For sufficiently small fixed $\xi$, every stop has actual last
exponent $b$ and actual ratio satisfying
\begin{equation}\label{eq:stops-actual-window}
 b_0-o(1)\le b\le b_1,
 \qquad 2.04\le r_d/b\le2.18.
\end{equation}
Here and throughout this section, the eventual bounds may depend on all
the fixed constants, including $\xi$.  The tag prime lies in a search
bin, whose exponent tends to infinity.  Thus it is different from the
last prime $P=w^b$ at every stop, once $z$ is sufficiently large.

\subsection{The average over a large prime bin}

At a stop, we will fix all prime factors except the tag prime.  The
arithmetic estimate below first lets that prime range over its whole
bin, imposing the aligned hit class even where the actual class does
not align.  Its lower average and absolute individual upper bound will
then permit removal of the small nonaligning subset.

For a prime $P$, write
\[
 V_{\rm all}(P)=\prod_{t<P}\left(1-\frac1t\right),
\]
where $t$ ranges over primes.

\begin{lemma}[Average for an aligned progression family]
\label{lem:stops-aligned-average}
Fix $C_s$.  Take $b_0$ sufficiently large in terms of $C_s$, and let
$b_1>b_0$ be fixed.  For every sufficiently small fixed $\xi>0$, the
following holds for all sufficiently large $z$.
Let $A/D$ be eligible in the sense of \eqref{eq:stops-eligible}, let
$P=w^b$ be a prime with $b\in[b_0-1,b_1]$, and let $q'$ be a squarefree
product of primes at most $z$ whose least prime factor is $P$ and all of whose prime factors
align to $A/D$.  Let $\mathcal P$ be a search bin, none of whose primes
divides $q'$.  Suppose that, for every $p\in\mathcal P$,
\begin{equation}\label{eq:stops-length-exponents}
 J_p=\frac{Y}{pq'},\qquad 2.03\le\log_PJ_p\le2.19.
\end{equation}
Define $T_p$ to count the integers $1\le n\le Y$ satisfying
$Dn\equiv A\pmod{pq'}$ and avoiding the prescribed classes at every
prime $t<P$.  Then
\begin{equation}\label{eq:stops-all-primes-lower}
 \sum_{p\in\mathcal P}T_p
 \ge .39V_{\rm all}(P)\sum_{p\in\mathcal P}J_p.
\end{equation}
There is also an absolute constant $C_U$, independent of
$C_s,b_0,b_1$ and of sufficiently small fixed $\xi$, such that
\begin{equation}\label{eq:stops-individual-upper}
 T_p\le C_UJ_pV_{\rm all}(P)\qquad(p\in\mathcal P).
\end{equation}
Both estimates are uniform in the eligible rational, the bin, the
product $q'$, and the prescribed residue classes.
\end{lemma}

The lower-bound proof writes $Dn-A=pq'm$ and counts integers
$m=\pm d_0\ell'$ with $d_0\le P^{.95}$ and $\ell'>P$ prime.
The prime factors of $d_0$ are chosen among the nonaligned primes
less than $P$ that do not divide $D$.  For such a prime $t$, divisibility
of $d_0$ by $t$ makes the avoidance condition automatic; otherwise
one of the $t-1$ unit classes for the varying prime is forbidden.
We isolate the resulting estimate before proving the arithmetic lemma.
For any subset $E$ of the primes less than $P$, put
\begin{equation}\label{eq:stops-euler-weight}
 W_E(d_0)=\prod_{\substack{t\in E\\t\nmid d_0}}
                    \left(1-\frac1{t-1}\right).
\end{equation}
The next lemma is uniform in $E$, including when $2\in E$.

\begin{lemma}\label{lem:stops-euler-probability}
Uniformly over all subsets $E$ of the primes less than $P$,
\begin{equation}\label{eq:stops-euler-lower}
 \sum_{\substack{d_0\le P^{.95}\\t\mid d_0\Rightarrow t\in E}}
       \frac{W_E(d_0)}{d_0}
 \ge e^{-\gamma}(.94-o(1)).
\end{equation}
\end{lemma}

\begin{proof}
For $t\in E$, let $V_t$ be independent nonnegative integer-valued
random variables with probabilities
\[
 \Pbb(V_t=0)=\frac{t-2}{t-1},\qquad
 \Pbb(V_t=j)=t^{-j}\quad(j\ge1).
\]
These masses sum to one, and the probability that
$\prod_{t\in E}t^{V_t}=d_0$ is exactly $W_E(d_0)/d_0$.  This
interpretation includes $t=2$: the mass at zero is then zero and
$\sum_{j\ge1}2^{-j}=1$.  In particular the zero factors in
\eqref{eq:stops-euler-weight} cause no exception to the probability
interpretation.

Let $U_t$ instead have the ordinary geometric probabilities
$\Pbb(U_t=j)=(1-1/t)t^{-j}$ for $j\ge0$.  For $j\ge1$,
\[
 \Pbb(V_t\ge j)=\frac{t^{1-j}}{t-1}
 \ge t^{-j}=\Pbb(U_t\ge j).
\]
Couple $U_t,V_t$ monotonically, independently over $t$.  Their mean
difference is
\[
 \E(V_t-U_t)
 =\frac{t}{(t-1)^2}-\frac1{t-1}
 =\frac1{(t-1)^2}.
\]
It follows that the excess log-product
$Z_E=\sum_{t\in E}(V_t-U_t)\log t$ has uniformly bounded mean,
since $\sum_t\log t/(t-1)^2<\infty$.  Extend the independent ordinary
variables to all primes $t<P$.  Then
\[
 \prod_{t\in E}t^{U_t}\le\prod_{t<P}t^{U_t}.
\]
The latter product has mass $V_{\rm all}(P)/n$ at every integer $n$
composed of primes below $P$.  Every $n\le P^{.94}$ has this property,
so
\[
 \Pbb\left(\prod_{t<P}t^{U_t}\le P^{.94}\right)
 =V_{\rm all}(P)\sum_{n\le P^{.94}}\frac1n
 =e^{-\gamma}(.94+o(1)).
\]
By Markov's inequality,
$\Pbb(Z_E>.01\log P)=O(1/\log P)$, uniformly in $E$.  Outside this
exception the preceding event implies
$\prod_{t\in E}t^{V_t}\le P^{.95}$.  Subtracting the exceptional
probability proves \eqref{eq:stops-euler-lower}; no independence
between these last two events is required.
\end{proof}

\begin{proof}[Proof of \cref{lem:stops-aligned-average}]
Every prime in the search bin eventually exceeds both $D$ and $P$.
The alignment of the prime factors of $q'$ implies $(D,q')=1$;
hence $(D,pq')=1$ for every $p\in\mathcal P$.  The hypothetical hit
condition and integrality of $n$ are therefore equivalent to
\begin{equation}\label{eq:stops-m-coordinate}
 Dn-A=pq'm,
 \qquad
 1\le\frac{A+pq'm}{D}\le Y,
 \qquad pq'm\equiv-A\pmod D.
\end{equation}
We count a subfamily of these integers in which $m$ has one prime
factor exceeding $P$.  Averaging over $p$ then imposes the remaining
avoidance conditions.

\paragraph{A common interval for the auxiliary integer.}
Partition $\mathcal P$ into
$\lceil w^{2C_s+10}\rceil$ equal-length subbins in prime coordinates.
In one such subbin let $R'$ be its lower endpoint, let $V$ be its
length, and put
\[
 J^\circ=\frac{Y}{R'q'},\qquad s^\circ=\log_P J^\circ.
\]
Uniformly over these subbins,
\begin{equation}\label{eq:stops-subbin-size}
 \frac{V}{R'}\asymp\xi w^{-2C_s-10},
 \qquad J_p=J^\circ(1+o(1)).
\end{equation}
Let $I$ be the intersection, over the real values of $p$ in the
subbin, of the real $m$-intervals specified by the middle condition
in \eqref{eq:stops-m-coordinate}.  The endpoints of one such interval
are
\[
 \frac{D-A}{pq'},\qquad \frac{DY-A}{pq'}.
\]
Relative to $DJ^\circ$, the variation of either endpoint is at most
\[
 O\left(\xi w^{-2C_s-10}
       \left(1+\frac{|A|}{DY}\right)\right)=o(1),
\]
by \eqref{eq:stops-eligible}.  Consequently
\begin{equation}\label{eq:stops-common-interval}
 |I|=(1-o(1))DJ^\circ.
\end{equation}
The height of this interval satisfies
\[
 \sup_{m\in I}|m|
 \le (D+|A|/Y)J^\circ(1+o(1))
 \le P^{s^\circ+.02},
\]
provided $b_0$ is sufficiently large in terms of $C_s$.  Indeed the
extra factor is bounded by a constant times $w^{C_s+2}$, which is
absorbed by $P^{.02}$ after such a choice.  All these statements are
uniform in the aligned progression family, its eligible rational,
and the subbin.

Among primes $t<P$, distinguish those aligning to $A/D$ from those
not aligning.  Let $E$ be the set of nonaligned primes $t<P$ that do
not divide $D$.  We count only integers of the form
\begin{equation}\label{eq:stops-factorization}
 m=\pm d_0\ell',\qquad
 d_0\le P^{.95},\qquad
 t\mid d_0\Longrightarrow t\in E,\qquad
 \ell'>P\ \hbox{prime},
\end{equation}
with $m\in I$.  Such a representation is unique: its factor $d_0$
is less than $P$, while the distinguished prime factor is greater than
$P$.  Furthermore, $(m,D)=1$, since $d_0$ has no prime factor dividing
$D$, and $D<P<\ell'$.  If $t<P$ aligns to $A/D$, then $t$ divides
neither $d_0$ nor $\ell'$, so $t\nmid m$, exactly as required to
avoid the class at this aligned prime.

For each permitted $d_0$, the number of integers
\eqref{eq:stops-factorization} in $I$, counting both signs, is at least
\begin{equation}\label{eq:stops-prime-factor-count}
 (1-o(1))\frac{|I|}{d_0(s^\circ+.02)\log P}.
\end{equation}
Here the error tends to zero uniformly in $d_0$, $E$, and all the
other data.  To prove this, split $I$ into its positive and negative
absolute-value pieces.  Remove $|m|<P^{-.02}|I|$, and discard any
remaining sign piece shorter than $P^{-.02}|I|$.  The total removed
length is $O(P^{-.02}|I|)$.  After division by $d_0$, every retained
interval has lower endpoint at least
\[
 \frac{P^{-.02}|I|}{d_0}
 \ge(1-o(1))D P^{s^\circ-.97}>P^{1.04},
\]
using \eqref{eq:stops-length-exponents} and
\eqref{eq:stops-common-interval}.  Its length divided by its upper
endpoint is at least $(1-o(1))DP^{-.04}$, and hence at least
$P^{-.05}$.  If its upper endpoint is $v$, then $v>P^{1.04}$ implies
\[
 vP^{-.05}\ge v^{9/10}.
\]
The uniform short-interval prime number theorem in
\cref{lem:prime-inputs} therefore applies.  Every upper endpoint has
logarithm at most $(s^\circ+.02)\log P$.  Summing the resulting
lower bounds for the retained sign pieces proves
\eqref{eq:stops-prime-factor-count}.  In particular their prime factors
automatically exceed $P$.

\paragraph{Sieving the varying prime.}

Fix one of the integers $m$ just counted.  The last congruence in
\eqref{eq:stops-m-coordinate} prescribes a single unit class for $p$
modulo $D$.  This uses $(A,D)=(q'm,D)=1$.  For each prime $t\mid D$,
increasing $p$ by $D$ changes the integer $n$ by $q'm$, which is a
unit modulo $t$.  Thus avoiding $a_t$ removes exactly one of the $t$
lifts of this base class modulo $Dt$.  Taken together, the remaining
classes modulo
\[
 Q=D\rad(D)
\]
form a set $\mathcal A$ of
$\varphi(\rad(D))$ unit classes.  Since
$\varphi(Q)=\varphi(D)\rad(D)$, their proportion among all unit
classes modulo $Q$ is exactly
\begin{equation}\label{eq:stops-lift-fraction}
 \frac{|\mathcal A|}{\varphi(Q)}=\frac1D.
\end{equation}
For $D=1$ this means the single unrestricted class and the same formula
holds.

For $t\in E$, avoiding the original class is equivalent to
\[
 pq'm\not\equiv Da_t-A\pmod t,
\]
with nonzero right side.  If $t\mid d_0$, this is automatic.  If
$t\nmid d_0$, then $m$ is a unit modulo $t$ because $\ell'>P$, and
exactly one nonzero $p$-class modulo $t$ is forbidden.  Thus the
product of the remaining survival factors is
$W_E(d_0)$ from \eqref{eq:stops-euler-weight}.
Let $N_{\mathcal I}$ be the number of primes in the present subbin
$\mathcal I$.  We claim that for any fixed $\gamma_1>0$, for all
sufficiently large $z$,
\begin{equation}\label{eq:stops-tag-prime-count}
 \#\{p\in\mathcal I:\ p\ \hbox{gives a local survivor for this }m\}
 \ge(1-\gamma_1)\frac{N_{\mathcal I}}D W_E(d_0),
\end{equation}
uniformly in all the data.  If the product vanishes, the claim is
trivial.  In every other case, the restricted primes are at least $3$;
their sieve densities $g(t)=1/(t-1)$ satisfy the hypotheses of
\cref{lem:fundamental}.

Here are the required progression and remainder bounds.  Choose a fixed
$A_2$ large enough, in terms of $\gamma_1$, to make the relative
fundamental-lemma error small, and use sieve level
$S=P^{A_2}$.  For a squarefree intersection modulus $l\le S$ on
the restricted primes and one class in $\mathcal A$, the forbidden
conditions specify a single reduced class modulo $Ql$, since
$(l,Q)=1$.  Siegel--Walfisz from \cref{lem:prime-inputs} gives its
count as
\[
 \frac{\li(R'+V)-\li(R')}{\varphi(Q)\varphi(l)}
 +O_{M_2}\left(\frac{R'}{(\log R')^{M_2}}\right).
\]
The exponent $M_2$ can be any sufficiently large fixed constant.  To
check the modulus range, observe that
\begin{equation}\label{eq:stops-sw-size-bounds}
 Q\le w^{2C_s},\qquad |\mathcal A|\le w^{C_s},\qquad
 S\le w^{b_1A_2},\qquad \log R'\ge w^{1/4}\log w.
\end{equation}
Thus $Ql$, the level, and the lift count are each bounded by fixed
powers of $\log R'$.  The main term, for each base class, has the
required multiplicative intersection density
$1/\varphi(l)=\prod_{t\mid l}1/(t-1)$.

For completeness, let $C$ be the fixed divisor-weight exponent in the
remainder of \cref{lem:fundamental}.  Summing over lifts and using the
crude bound $\sum_{l\le S}\tau(l)^C\ll S^{C+1}$, the total remainder
is at most
\[
 O_{M_2}\left(
 |\mathcal A|S^{C+1}\frac{R'}{(\log R')^{M_2}}
 \right).
\]
In a nonzero case, uniformly over subsets of restricted primes,
\[
 W_E(d_0)\gg\frac1{\log P}.
\]
Indeed the product over all $3\le t<P$ is of this order, by Mertens'
theorem and the convergent positive product in the identity
\[
 1-\frac1{t-1}
 =\left(1-\frac1t\right)
   \left(1-\frac1{(t-1)^2}\right).
\]
Also $N_{\mathcal I}\asymp V/\log R'$, uniformly by the same
prime number theorem, or by Siegel--Walfisz with modulus one.  Write
$\kappa=2C_s+10$.  The ratio of the displayed remainder to
$(N_{\mathcal I}/D)W_E(d_0)$ is consequently at most
\[
 O_{M_2,\xi}\left(
 \frac{D|\mathcal A|w^\kappa S^{C+1}\log P}
      {(\log R')^{M_2-1}}
 \right)
 \ll_{M_2,\xi}
 \frac{w^{4C_s+10+(C+1)b_1A_2}\log P}
      {(\log R')^{M_2-1}}=o(1)
\]
on taking $M_2$ sufficiently large in terms of the fixed constants.
This also justifies subtracting the two endpoint estimates for the
short subbin.  The logarithmic-integral main term can be replaced by
$N_{\mathcal I}$ with relative error $o(1)$.  Finally use
\eqref{eq:stops-lift-fraction}, and take the fundamental-lemma relative
error sufficiently small compared with $\gamma_1$.  This proves
\eqref{eq:stops-tag-prime-count}.  Its error bound is uniform in $m$;
large values of $m$ change only reduced residue classes.

\paragraph{Summing the lower bound and bounding an individual count.}
For each $d_0$, first count the integers $m$ using
\eqref{eq:stops-prime-factor-count}, and for each such $m$ count its
admissible varying primes using \eqref{eq:stops-tag-prime-count}.  A pair
$(p,m)$ determines a unique integer $n$ by
\eqref{eq:stops-m-coordinate}.  Unique factorization in
\eqref{eq:stops-factorization} prevents duplicate counting for fixed
$p$.  Hence \cref{lem:stops-euler-probability} and
\eqref{eq:stops-common-interval} give, in one subbin,
\[
 \sum_{p\in\mathcal I}T_p
 \ge (1-\gamma_1-o(1))
       \frac{N_{\mathcal I}J^\circ}{(s^\circ+.02)\log P}
       e^{-\gamma}(.94-o(1)).
\]
Mertens' theorem gives
$V_{\rm all}(P)\sim e^{-\gamma}/\log P$.  We have
$s^\circ+.02\le2.22$ eventually, and
\[
 \frac{.94}{2.22}>.423.
\]
Choose $\gamma_1$ as a sufficiently small absolute constant; for
example $\gamma_1=.05$ leaves a limiting coefficient greater than
$.402$.  Since $J_p=J^\circ(1+o(1))$ uniformly within a subbin,
summing over all subbins proves \eqref{eq:stops-all-primes-lower}.
All relative errors used here are uniform, so the growing number of
subbins and the number of pairs $(d_0,m)$ do not accumulate additional
errors.

For the upper bound, the hit conditions give one ordinary progression
of step $pq'$ and coordinate length $J_p+O(1)$.
Every sieve prime below $P$ is
coprime to this step.  Apply the upper form of
\cref{lem:fundamental} using primes up to $P^\theta$, for a
sufficiently small fixed absolute $\theta>0$, and a level such as
$P^{1/2}$.  The relative sieve parameter is a sufficiently large
absolute constant; the summed $O(1)$ intersection remainders are
$P^{1/2+o(1)}$, whereas $J_p\ge P^2$.  Mertens' theorem gives
\[
 \prod_{t<P^\theta}(1-1/t)\ll_\theta V_{\rm all}(P).
\]
This proves \eqref{eq:stops-individual-upper} with an absolute
$C_U$, independent of $C_s,b_0,b_1$ and of sufficiently small $\xi$.

\end{proof}

\subsection{Comparison at the actual stops}

The preceding lemma averages over every prime in a bin.  To apply it
to actual stops, we must show that fixing the other prime factors
leaves the entire owner-aligning subset available.  This is the role
of the representative tests in the stop rule.

\begin{proposition}[Comparison at stops]
\label{prop:stopped-comparison}
There is an absolute constant $\eta_0>0$ with the following property.
Fix $C_s$ and the regular-box parameters.  Take $0<\eta\le\eta_0$,
then take $b_0$ sufficiently large in terms of $C_s$, and take any
fixed $b_1>b_0$.  For all sufficiently small fixed $\xi>0$ and all
sufficiently large $z$, the stops of \cref{def:stops-rule} satisfy
\begin{equation}\label{eq:stops-comparison}
 \sum_{d\ \mathrm{stopped}} S_d(x(P_d))
 \ge
 \sum_{d\ \mathrm{stopped}}\mu_d P_\ee(r_d,x(P_d)),
\end{equation}
where $P_d$ denotes the last prime of the tuple and both cutoffs are
strict below $P_d$.
\end{proposition}

\begin{proof}
We first explain exactly which prime is averaged.  Partition the stopped
nodes by their ordered bin labels, their tag position, and all prime
choices other than the tag prime $p$.  In a nonempty group, let $A/D$
be the owner of its tag bin.  Then $p$ ranges over the entire subset of
that bin aligning to $A/D$.

To justify this assertion, replace $p$ by any other aligning prime in
the bin.  All earlier choices remain fixed, so no earlier tag appears.
At the designated position the new prime aligns to the same owner;
hence the tag position and tag rational remain unchanged.  Every
subsequent all-prefix alignment test has the same truth value, since
the other primes are fixed and the replacement aligns.  The
representative tests, lengths, and bin multiplicities depend only on
the labels.  Admission is preserved by \eqref{eq:stops-safety}, and
the tag-bin multiplicity one prevents a change in prime ordering or a
collision with another factor.  These observations apply to each
earlier prefix as well, so no earlier stop is introduced.  This proves
the assertion about the entire aligning subset.

In this group put $q'=d/p$.  The last prime $P$ is a factor of $q'$
and is fixed.  All primes of $q'$ are at least $P$, and they align to
$A/D$.  First allow \emph{all} primes $p$ in the tag bin, assigning
to $p$ the hypothetical hit class aligned to $A/D$.  These hypothetical
choices need not be actual stops; their local counts are auxiliary.
Let $\mathcal P$ denote this tag bin, and use the hypothetical counts
$T_p$ of \cref{lem:stops-aligned-average}.  On its aligning subset,
$T_p=S_{pq'}(x(P))$.

For all choices of $p$ in the tag bin, put $b=x(P)$ and use the gap
$r_{pq'}=\log_w(Y/(pq'))-a_\star+2$.  Thus
\[
 \log_PJ_p=\frac{r_{pq'}}b+\frac{a_\star-2}{b}.
\]
The bin labels and representative safety tests are fixed, so the
movement bound applies also to these hypothetical choices.  By
\eqref{eq:stops-actual-window}, after enlarging $b_0$ if necessary,
all of them satisfy \eqref{eq:stops-length-exponents}.  The search
bin is disjoint from the prime factors of $q'$ because the tag bin
has multiplicity one.  All hypotheses of
\cref{lem:stops-aligned-average} therefore hold.

For sufficiently large $z$, the actual cutoff lies in the fixed
interval $[b_0-1,b_1]$.
Applying \cref{lem:local-reference} on this interval, together with the
product identity $V_{\rm all}(P)=V_0V(b)$, gives
\[
 \frac{\mu_{pq'}P_\ee(r_{pq'},b)}{J_pV_{\rm all}(P)}
 \le f(r_{pq'}/b)+O(1/b)+o(1).
\]
The constants are uniform for the fixed enlarged window $[b_0-1,b_1]$ and for
the ratio window in \eqref{eq:stops-actual-window}.  On that ratio
window, the near-$2$ formula and monotonicity of $f$ give
\[
 f(r_{pq'}/b)\le f(2.18)
 <\frac{2(1.95)(.166)}{2.18}<.30.
\]
Here $e^\gamma<1.95$ follows, for example, from
$\gamma\le\sum_{j=1}^6j^{-1}-\log6<.66$, and
$\log(1.18)<.166$ follows from
$\log(1+u)\le u-u^2/2+u^3/3$ at $u=.18$.
Thus first choosing $b_0$ large and subsequently taking $z$ large gives
\begin{equation}\label{eq:stops-reference-upper}
 \mu_{pq'}P_\ee(r_{pq'},b)\le .33 J_pV_{\rm all}(P)
\end{equation}
uniformly over the group and over its hypothetical choices.

Apply \eqref{eq:stops-all-primes-lower} to the full tag bin.
Let $\mathcal G$ be the aligning subset in the tag bin.  Its complement
contains at most an $\eta$ fraction of the bin primes.  Since the
$J_p$'s vary by at most $1+\xi$ within the whole bin, for $\xi\le1$
we have
\[
 \sum_{p\notin\mathcal G}T_p
 \le 2C_U\eta V_{\rm all}(P)
                    \sum_{p\ \mathrm{in\ bin}}J_p.
\]
Fix an absolute $\eta_0<1/2$ so small that
$2C_U\eta_0<.05$.  For $0<\eta\le\eta_0$, subtracting this bound
from \eqref{eq:stops-all-primes-lower} and using
\eqref{eq:stops-reference-upper} gives
\[
 \begin{split}
 \sum_{p\in\mathcal G}S_{pq'}(x(P))
 &=\sum_{p\in\mathcal G}T_p\\
 &\ge .34V_{\rm all}(P)\sum_{p\ \mathrm{in\ bin}}J_p\\
 &\ge\sum_{p\in\mathcal G}\mu_{pq'}P_\ee(r_{pq'},b).
 \end{split}
\]
The last inequality uses $.34>.33$ and positivity of $J_p$; it does
not require a distribution hypothesis for the survivors inside the
owner subset.  Every stopped node belongs to exactly one of the
groups considered.  Summing the group inequalities proves
\cref{prop:stopped-comparison}.

\end{proof}

\subsection{Almost every hard endpoint lies below a stop}

We next explain why the same stops capture all but an arbitrarily small
scaled harmonic mass of the hard compact endpoints.  In what follows,
the harmonic mass of a family of tuples is $\sum_d1/d$.

\begin{proposition}[Stopping hard endpoints]
\label{prop:hard-stopping}
Fix a compact gap bound $K$, the regular-box parameters, the bounded
candidate lists from \cref{prop:inverse}, and the exponent $C_s$
from \cref{cor:hard-endpoints}.  Fix
$0<\eta\le\eta_0$ as in \cref{prop:stopped-comparison}.
For every $\delta>0$, take $b_0$ sufficiently large in terms of $C_s$
and $K$, then take $b_1/b_0$ sufficiently large in terms of the desired
error, and finally take $\xi$ sufficiently small.  For these fixed
choices, the regular included hard endpoints with $r_d\le K$ having
no stopped prefix have harmonic mass at most
\begin{equation}\label{eq:stops-hard-mass}
 (\delta+o(1))B^{-2}.
\end{equation}
The eventual $o(1)$ may depend on all the fixed choices.  In particular,
$b_1$ and $\xi$ are fixed before $z$ tends to infinity.
\end{proposition}

\begin{proof}
For a hard endpoint, \cref{cor:hard-endpoints} gives a rational on one
of its candidate lists that satisfies \eqref{eq:stops-eligible} and
aligns to every path factor of exponent greater than $C_s$.  We need
two facts about this path: its tag must be this rational, and
it must have an earlier even visit in the stop window.  We establish
the tag statement in the full-box product measure, and the visit
statement in the unconditioned path process.  Their exceptional masses
will then be added.

\paragraph{Identifying the tag.}
Except for $o(B^{-2})$ mass, the path acquires this same rational as
its tag during its search positions, as we now show.

Fix one regular full candidate box and one eligible rational $A/D$
on its list.  Each of its search bins has multiplicity one, and
\cref{prop:regular-boxes} supplies at least $c_7\log w$ such bins.
For clarity, the count follows because all nonsearch factors together
decrease the gap by at most
\[
 C_{\rm len}\log B\,(w^{1/4}+h)=o(w^{1/2}).
\]
The search factors come first and must therefore decrease the gap from
$r_0\asymp B\asymp w\log w$ to at most $w^{1/2}$.  A standard
transition decreases the gap by at most a bounded multiplicative
factor, including at the initial extended start.  At least a constant
multiple of $\log w$ search factors is necessary.

At a hard endpoint for the chosen rational, all these search primes
align to it.  If none of their bins is near-full for this rational,
each of the bins has more than an $\eta$ fraction of nonaligning
primes by count.  For primes in $[R,(1+\xi)R]$, this is a harmonic
fraction at least $\eta/(1+\xi)$.  The search choices are independent
in the full product box, so the fraction of its harmonic measure with
all those alignments is at most
\[
 \left(1-\frac{\eta}{1+\xi}\right)^{c_7\log w}=o(1).
\]
This bounds the relevant hard event in the full box without asserting
independence after conditioning on hardness.

Otherwise some visited search bin is near-full for this rational,
which is therefore its owner.  The hard endpoint's choice in this bin
aligns to its owner, so a tag has been acquired by that position.  If
the first tag is a different rational, the prime at the earlier tag
position aligns both to $A/D$ and to the distinct eligible owner of
that bin.  By \eqref{eq:stops-overlap}, its possible choices form a
uniformly negligible fraction, even after summation over
$C_{\rm len}\log B$ positions.  Indeed in a search bin this fraction
is at most a constant, depending on fixed $\xi$, times
\[
 \frac{L+\log R}{R},
\]
and $R\ge\exp(w^{1/4}\log w)$.

There are at most the fixed list bound times the fixed candidate count
such rational choices per box.  The union of the preceding exceptions
therefore costs $o(1)$ of the full candidate-box harmonic measure.
The total mass of the enlarged boxes is $O_K(B^{-2})$, independently
of sufficiently small $\xi$, by \cref{prop:regular-boxes}.  Summing
over the boundedly many candidates gives a total $o(B^{-2})$
exception.  Outside it, a hard endpoint has the correct tag by the end
of its search positions.

\paragraph{Finding an even visit before the compact endpoint.}
Apply
\cref{lem:marked-visits} to the continuous standard path.  For any
prescribed probability error $\varepsilon_1>0$, with $b_0$ fixed,
choosing $b_1/b_0$ sufficiently large ensures, uniformly from the
allowed starting ratios as the initial gap tends to infinity, an even
arrival with ratio in $[2.08,2.14]$ and last exponent in
$[2b_0,b_1/2]$, except with probability at most
$\varepsilon_1+o(1)$.  The proof of that lemma marks the first draw
from an odd regeneration when its ratio lies in $[2.09,2.13]$, using
a pre-draw gap in $[10b_0,b_1]$.  In particular the arrival has the
stronger interior margins
\begin{equation}\label{eq:stops-marked-slack}
 \frac{10b_0}{3.13}\le x\le\frac{b_1}{3.09},
 \qquad 2.09\le s\le2.13.
\end{equation}
The choice of $b_1$ depends on the probability budget and fixed $b_0$,
not on $z$ or on $\xi$.

Use \cref{prop:prime-coupling} for standard prime paths with lower
exponent cutoff $1$.  Actual included paths are a subset of these
standard paths; their rules agree above exponent $2$.  High-ratio
histories ending in the fixed compact window have negligible scaled
mass by \cref{lem:high-states}.  On the remaining histories the
coupling matches states and exits, with a failure probability tending
to zero.  The interior margins in \eqref{eq:stops-marked-slack}
therefore give the asserted slightly wider visit windows on matched
prime paths.

This visit lies in the segment being compared whenever the matched
prime history ends in the compact window.  In fact its arrival gap is
at least a fixed multiple of $b_0$.  Taking $b_0$ sufficiently large
relative to $K$, that gap exceeds the matched compact endpoint gap.
Since gaps decrease along the path, a raw continuous marked visit
cannot occur after this endpoint.  Matching exits puts the visit
before any cutoff exit as well.  Thus the coupling cannot lose the
needed visit merely by stopping the compared segment at the compact
endpoint.

The compact-history event estimate in \cref{prop:prime-occupation}
now bounds the harmonic mass of compact endpoints missing such a visit
by
\[
 C_K(\varepsilon_1+o(1))B^{-2}+o(B^{-2}),
\]
where $C_K$ depends on the fixed compact window and the starting bounds,
not on the later bin width.  Choose $\varepsilon_1$ sufficiently small
in terms of $\delta/C_K$, and choose $b_1$ accordingly.

\paragraph{Verifying the stop conditions.}
Consider a regular hard endpoint with the correct tag and with this
even visit.  Every factor up to the visit has exponent at least the
visit's last exponent, which exceeds $C_s$ by the choice of $b_0$.
Hence all these factors align to its tag rational.  Its search
positions have already passed: their lower exponents are at least
$w^{1/4}$, whereas the visit exponent is bounded by fixed $b_1$.
The tag is therefore attained by the visit, and its bin has multiplicity
one.  The visit-prefix length is at most the regular endpoint length
$C_{\rm len}\log B$.

By regularity, every earlier nonempty even prefix, including the visit,
has actual margin greater than $10\Delta$ in the admission test.
The representative movement is at most $\Delta$, and changing the
last exponent changes $\mathcal H$ by at most $2h$.  For sufficiently large
$z$, the representative margin is consequently at least
$10\Delta-\Delta-2h>3\Delta$.  Thus
\eqref{eq:stops-safety} holds.  Finally the ratio and exponent slack
in the visit windows, the prefix movement bound, and a sufficiently
small fixed $\xi$ imply \eqref{eq:stops-window}.  All stop conditions
hold at this visit.  The branch has therefore stopped there or earlier.

Combining the $o(B^{-2})$ incorrect-tag exception with the arbitrarily
small missing-visit mass proves \eqref{eq:stops-hard-mass}.
\end{proof}

The order of choices in these two propositions is worth making
explicit.  The constant $C_U$ in
\eqref{eq:stops-individual-upper} is absolute, so the owner threshold
$\eta$ can be fixed before either stop-window endpoint.  The choice
of $b_0$ absorbs $C_s$, the local reference error, the interval-height
margin, and the compact gap range.  Only the ratio $b_1/b_0$ must then
be enlarged to make the marked-visit error as small as desired.
Afterwards $\xi$ is chosen for representative movements and all earlier
freezing conditions.  Every progression modulus and prime-sieve level
used here remains a fixed power of $\log R'$ for these choices, however
large the fixed $b_1$ is.  The corresponding asymptotic thresholds may
therefore wait until this final fixed $\xi$ has been chosen.
